
\section{Methodology}
\subsection{Overall Architecture}
DHC follows the edge-batch procedure in Algorithm~\ref{alg:dhc_forward} 
The model first predicts positive and negative links with the current states. 
It then updates states after the prediction step. Negative samples are used for the loss only and do not
update memories. The representation has two complementary branches. 
The relation-aware context branch uses text-aware node initialization, dynamic relation states, and recent-neighbor subgraphs
to explore relation semantics and local heterogeneous context. 
The sequential event-memory branch stores compact chronological interaction history to preserve event order and repeated interaction signals. 
Then, a query-level gate combines the two branches before link scoring. 
The overall architecture of our framework is illustrated in Fig.~\ref{fig:dhc_architecture}. 
\begin{figure}[t]
\centering
\includegraphics[width=\linewidth]{figures/2-framework.png}
\caption{ 
Framework of DHC which combines text-aware initialization, relation-aware temporal messaging, dynamic node/relation memory, recent-neighbor subgraph encoding.
and event-memory fusion for temporal link prediction on heterogeneous graphs.
}
\label{fig:dhc_architecture}
\vspace{-2mm}
\end{figure}

\begin{algorithm}[t]
\caption{DHC edge-batch forward and chronological state update.}
\label{alg:dhc_forward}
\small
\begin{algorithmic}[1]
\Require Chronological batch $B=\{(u_i,r_i,v_i,t_i,\mathbf{e}_i,y_i)\}$, current states $\mathbf{M},\mathbf{R}$, recent-neighbor store $\mathcal{N}$, sequential event memory $\mathcal{E}$, and sampled negatives $v_i^{-}$.
\Ensure Positive and negative logits, target-edge loss, and updated dynamic states.
\State Build augmented edge features $\tilde{\mathbf{e}}_i=[\mathbf{e}_i\Vert\psi(r_i)]$ when edge-type augmentation is enabled.
\State Collect query nodes from $\{u_i,v_i,v_i^{-}\}$.
\State Encode query nodes with the relation-aware context branch using base node features, relation states, and recent-neighbor subgraphs.
\State Read query-node states from the sequential event-memory branch.
\State Fuse the two branch outputs with a vector gate.
\State Compute $s_i^{+}=s_\theta(u_i,r_i,v_i,t_i)$ and $s_i^{-}=s_\theta(u_i,r_i,v_i^{-},t_i)$.
\State Compute BCE loss only on target edges with $y_i=1$.
\State Generate directional relation-time messages from the observed edges in $B$.
\State After back-propagation, update node memory with the latest node messages in the batch.
\State Insert observed edges into the recent-neighbor store.
\State Update sequential event memory with observed edges and augmented edge features.
\State Update relation memory with the latest relation messages in the batch.
\end{algorithmic}
\end{algorithm}

\subsection{Input Encoding and Relation-time Messages}
For node $u$, DHC builds a base representation from its text-derived feature vector $\mathbf{x}_u\in\mathbb{R}^{d_x}$, node type embedding $\mathbf{a}_{\tau_u}$, and current node memory:
\begin{equation}
\mathbf{h}^{0}_{u}(t)=f_{\mathrm{node}}\left(\mathbf{x}_{u}\Vert\mathbf{a}_{\tau_u}\Vert\mathbf{M}_{u}(t^{-})\right),
\label{eq:base_node}
\end{equation}
where $f_{\mathrm{node}}$ is a two-layer MLP. This step injects semantic attributes and type information before temporal message passing.
For an elapsed time $\Delta t$, DHC uses a sinusoidal encoding $\Phi(\Delta t)$ as below:
\begin{equation}
\Phi(\Delta t)=\left[\sin\left(\frac{\Delta t}{\omega_1}\right),\ldots,\sin\left(\frac{\Delta t}{\omega_q}\right),
\cos\left(\frac{\Delta t}{\omega_1}\right),\ldots,\cos\left(\frac{\Delta t}{\omega_q}\right)\right],
\label{eq:time_encoding}
\end{equation}
\rev{where $q=\lfloor d_t/2\rfloor$ and $\omega_1,\ldots,\omega_q$ are fixed log-spaced frequencies.}
If $d_t$ is odd, the last dimension is zero-padded. 
Each relation also has a learned embedding. If relation memory is enabled, the effective relation representation is:
\begin{equation}
\boldsymbol{\rho}_{r}(t)=\mathbf{a}_{r}+W_{\rho}\mathbf{R}_{r}(t^{-}).
\label{eq:relation_state}
\end{equation}
Thus, a relation is represented by both its type embedding and its current dynamic state. 
\rev{For an observed event $(u,r,v,t,\mathbf{e})$, the implementation optionally augments the edge feature with a relation one-hot vector, $\tilde{\mathbf{e}}=\mathbf{e}\Vert\psi(r)$.}
\rev{The raw message inputs are}
\begin{equation}
\mathbf{z}_{u\leftarrow v}=\mathbf{M}_{v}(t^{-})\Vert\mathbf{M}_{u}(t^{-})\Vert\tilde{\mathbf{e}}\Vert\Phi(t-t_{u}^{\mathrm{last}}),
\end{equation}
\rev{DHC then applies a relation-conditioned gate:}
\begin{equation}
\mathbf{m}_{u\leftarrow v}=W_m\mathbf{z}_{u\leftarrow v}\odot\sigma(W_g\boldsymbol{\rho}_{r}(t)) + W_b\boldsymbol{\rho}_{r}(t),
\label{eq:rel_time_msg}
\end{equation}
and the same form is used for $\mathbf{m}_{v\leftarrow u}$. The gate lets relation semantics modulate the temporal event message instead of treating the relation as a static label.

\subsection{Dynamic States and Recent-neighbor Storage}
DHC maintains node memory $\mathbf{M}_u$ and relation memory $\mathbf{R}_r$ with GRU cells. 
If a node or relation receives multiple messages in the same mini-batch, only the latest timestamped message is committed:
\begin{align}
\mathbf{M}_u(t) &= \mathrm{GRU}_n\left(W_n\mathbf{m}_{u}(t),\mathbf{M}_u(t^-)\right),\\
\mathbf{R}_r(t) &= \mathrm{GRU}_r\left(W_r\frac{\mathbf{m}_{u\leftarrow v}(t)+\mathbf{m}_{v\leftarrow u}(t)}{2},\mathbf{R}_r(t^-)\right).
\end{align}
Node memory captures endpoint-specific historical behavior, while relation memory
\rev{captures global non-stationarity of each relation type. The recent-neighbor store records a bounded interaction window for each node.}
When $(u,r,v,t,\tilde{\mathbf{e}})$ is observed, both directions are inserted with relation id, timestamp, event id, and augmented edge feature.  
For each node, DHC stores the most recent $K$ undirected interactions. Given a query node $u$, the stored events provide local temporal context:
\begin{equation}
\mathbf{c}_u(t)=\sum_{j\in\mathcal{N}_K(u,t)}\alpha_{uj}(t)\,W_v\left(\mathbf{M}_j(t^-)\,\Vert\,\boldsymbol{\rho}_{r_{uj}}(t)\,\Vert\,\phi(t-t_{uj})\right).
\end{equation}
The weights $\alpha_{uj}$ are computed by attention or replaced with uniform mean weights. This context helps the model use recent relation-aware evidence without scanning the full history.

\subsection{Dual-branch Query Representation}
For link prediction, the recent-neighbor store returns a small temporal subgraph around the query endpoints. 
Edge attributes in this subgraph combine relation representation, augmented edge features, and time encoding. 
Relation-aware branch produces the context representation $\mathbf{h}^{C}_u(t)$ by fusing the base node
representation and the local temporal subgraph representation:
\begin{equation}
\mathbf{h}^{C}_u(t)=f_N\left(\mathbf{h}^{0}_u(t)\,\Vert\,\mathbf{c}_u(t)\right).
\end{equation}
\rev{In parallel, the sequential event-memory branch maintains a compact chronological state $\mathbf E_u\in\mathbb R^{d_m}$ for every node, separate from DHC's relation-aware node memory. For an observed event $(u,v,t,\tilde{\mathbf e})$, it constructs two directional messages:}
\begin{equation}
\rev{\begin{aligned}
\mathbf q^{E}_{u\leftarrow v}(t)
&=\mathbf E_u(t^-)\Vert\mathbf E_v(t^-)\Vert\tilde{\mathbf e}
\Vert\Phi(t-t_u^{\mathrm{last},E}),\\
\mathbf q^{E}_{v\leftarrow u}(t)
&=\mathbf E_v(t^-)\Vert\mathbf E_u(t^-)\Vert\tilde{\mathbf e}
\Vert\Phi(t-t_v^{\mathrm{last},E}).
\end{aligned}}
\label{eq:event_memory_message}
\end{equation}
\rev{When several events address the same node in one chronological batch, the branch retains the message with the latest timestamp. Its state and query representation are therefore}
\begin{equation}
\rev{\begin{aligned}
\bar{\mathbf q}^{E}_{u}(t)
&=\operatorname{Last}\!\left(\{\mathbf q^{E}_{u\leftarrow j}(t_i):t_i\le t\}\right),\\
\mathbf E_u(t)&=\operatorname{GRU}_{E}\!\left(\bar{\mathbf q}^{E}_{u}(t),\mathbf E_u(t^-)\right),\\
\mathbf h^{E}_{u}(t)&=\mathbf E_u(t^-).
\end{aligned}}
\label{eq:event_memory_update}
\end{equation}
\rev{Equation~\eqref{eq:event_memory_update} is applied only when node $u$ receives an observed message; otherwise $\mathbf E_u(t)=\mathbf E_u(t^-)$. The pre-event state $\mathbf h^{E}_{u}(t)$ is used for scoring to prevent leakage. After scoring and optimization, Eqs.~\eqref{eq:event_memory_message}--\eqref{eq:event_memory_update} are applied only to observed target/context events; sampled negatives never update the state. This branch retains endpoint-specific event order and recurrence. In contrast, the context branch uses text/type features, DHC node memory, the recurrent global state of each relation type, and an explicitly materialized bounded recent-neighbor subgraph. Their parameters and state buffers are separate; they meet only in the following query-level gate.}
The two branches are fused by a vector gate:
\begin{equation}
\mathbf{g}_{u}(t)=\sigma\left(W_F\left(\mathbf{h}^{C}_{u}(t)\Vert\mathbf{h}^{E}_{u}(t)\right)+\mathbf{b}_F\right),
\end{equation}

\begin{equation}
\mathbf{h}_{u}(t)=\mathbf{g}_{u}(t)\odot\mathbf{h}^{E}_{u}(t)+\left(1-\mathbf{g}_{u}(t)\right)\odot\mathbf{h}^{C}_{u}(t).
\end{equation}

The gate allows each query node to balance local heterogeneous context and sequential event history.
\subsection{Link Scoring, Loss, and Negative Sampling}
For a candidate edge $(u,r,v,t)$, DHC scores:
\begin{equation}
s_{\theta}(u,r,v,t)=f_s\left(\mathbf{h}_{u}(t)\Vert\mathbf{h}_{v}(t)\Vert\boldsymbol{\rho}_{r}(t)\Vert\Phi(t-t_u^{\mathrm{last}})\right),
\label{eq:score}
\end{equation}
where $f_s$ is an MLP that outputs a logit. 
Training samples one negative destination $v_i^{-}$ for each observed event. The loss is applied only to supervised target edges in the batch:
\begin{equation}
\mathcal{L}_{\mathrm{link}}=\frac{1}{|B^{+}|}\sum_{i\in B^{+}}w_{r_i}\left[
\operatorname{BCE}(s_i^{+},1)+\operatorname{BCE}(s_i^{-},0)\right],
\label{eq:link_loss}
\end{equation}
where $B^{+}=\{i\in B\mid y_i=1\}$, $s_i^{+}=s_{\theta}(u_i,r_i,v_i,t_i)$, $s_i^{-}=s_{\theta}(u_i,r_i,v_i^{-},t_i)$, and $w_{r_i}$ is an optional relation-specific loss weight. Context edges with $y_i=0$ do not contribute to the loss, but they still update states after prediction.

The training sampler mixes three negative-destination sources:
\begin{equation}
P(v^{-}\mid v)=\rho_tP_{\mathrm{type}}(v^{-}\mid \tau_v)+\rho_pP_{\mathrm{pop}}(v^{-})+(1-\rho_t-\rho_p)P_{\mathrm{unif}}(v^{-}).
\end{equation}
The ratios $\rho_t$ and $\rho_p$ are tunable hyperparameters. 
They control the balance between type-consistent, popularity-biased, and uniform negative destinations. 
The default setting, $\rho_t=0.6$ and $\rho_p=0.2$, is an empirical trade-off. 
It makes most negatives type-compatible with the positive destination, adds a smaller fraction of hard popular negatives, and keeps the remaining samples uniformly drawn for diversity. 
Evaluation uses $K=20$ uniformly sampled negatives per positive target edge.
